% !PRINTED 504
\weberpage{504}

%<<CONT
trischen Gruppe zu einer Gruppe $Q$ gehört, so gehört sie auch innerhalb jeder anderen Gruppe $P$, von der $Q$ ein Theiler ist, zu $Q$.

Der zweite Theil aber folgt aus den Grundeigenschaften der \gesperrt{Galois'schen} Gruppe.

Um dies einzusehen, wollen wir folgende Bezeichnung einführen. Wird auf eine Function $\psi$ des Körpers $N$ irgend eine Permutation $\pi$ der Gruppe $P$ ausgeübt, so möge $\psi$ in $\psi(\pi)$ übergehen; wird auf $\psi(\pi)$ eine Permutation $\pi'$ ausgeübt, so ist der Erfolg derselbe, als ob auf $\psi$ die zusammengesetzte Permutation $\pi\pi'$ ausgeübt wäre, d.\,h. es ist
\weq{6}{\psi(\pi)(\pi') = \psi(\pi\pi').}
Bleibt nun $\psi$ ungeändert durch die Permutation $\varkappa$, so ist
\[
\psi = \psi(\varkappa),
\]
und hierin kann nach \S.\,149, a) jede Permutation $\pi$ aus der Gruppe $P$ angewandt werden. Dadurch ergiebt sich
\[
\psi(\pi) = \psi(\varkappa\pi).
\]
Wenn daher $\psi(\pi)$ auch gleich $\psi$ ist, so ist auch $\psi(\varkappa\pi) = \psi$, d.\,h. die Permutationen, die $\psi$ ungeändert lassen, bilden eine Gruppe, w.\,z.\,b.\,w.

Ist $\psi$ eine zur Gruppe $Q$ gehörige Function in $N$, ist also
\weq{7}{\psi(\varkappa) = \psi(\varkappa_1) = \cdots = \psi(\varkappa_{q-1}) = \psi,}
und $\pi_1$ irgend eine nicht zu $Q$ gehörige Permutation von $P$, so können wir $\pi_1$ auf die Gleichungen (7) anwenden und erhalten
\[
\psi(\varkappa\pi_1) = \psi(\varkappa_1\pi_1) = \cdots = \psi(\varkappa_{q-1}\pi_1) = \psi_1,
\]
und $\psi_1$ ist von $\psi$ verschieden. Wenn umgekehrt für irgend ein $\pi$
\weq{8}{\psi_1 = \psi(\pi) = \psi(\varkappa\pi_1)}
ist, so muss $\pi$ unter der Reihe der Substitutionen
\weq{9}{Q\pi_1 = \varkappa\pi_1,\ \varkappa_1\pi_1,\ldots \varkappa_{\nu-1}\pi_1}
% [sic] Der Druck setzt im Endindex \nu-1, obwohl Q in (2) und (7) bis q-1 indiziert wird; quellentreu beibehalten.
enthalten sein; denn durch Anwendung von $\pi^{-1}$ auf (8) folgt, dass $\varkappa\pi_1\pi^{-1}$ in $Q$ enthalten, etwa gleich $\varkappa_1$ sein muss; daraus aber folgt
\[
\pi = \varkappa_1^{-1}\varkappa\pi_1,
\]
was in der That in (9) enthalten ist.

Wir können demnach folgenden Satz aussprechen:

4. \gesperrt{Eine zu $Q$ gehörige Function $\psi$ geht durch alle Permutationen einer Nebengruppe $Q\pi$, und}
%CONT>>
